\originalchapter{模型进步的数据与计算条件}\label{ch:progress}
\originalsection{把进步拆成可检验的变化}\label{sec:progress-factors}
语言模型在多个任务上持续改善, 并不只因为参数数目增加. 参数决定可表示的函数类, 数据决定学习压力, 优化决定有限预算下找到什么参数, 系统实现决定多少计算实际可用, 后训练决定哪些行为被强化, 推理协议决定一次任务能生成与检查多少候选. 将这些因素分开, 才能提出可检验的问题.

例如要检验 ``更多数据导致更好的模型'', 应在同一架构、优化与评价协议下改变训练 token 预算, 同时记录训练是否重复旧数据. 要检验 ``清理数据更有效'', 应保持模型、训练计算和 token 数基本一致, 比较清理前后的资料. 要检验 ``推理能力提升'', 则应分别报告单次生成与多候选选择, 否则搜索预算的变化会被误认为权重的变化. 这些实验不能完全控制现实系统的每个因素, 但至少把因果主张变成可反驳的比较.

训练对序列位置可以并行, 是注意力模型易扩展的工程条件之一. 同一训练文档的所有真实前缀已经给定, 因而可一次计算大量 token 损失; 自回归生成仍有逐 token 的依赖. 硬件、精度、算子融合与数据搬运的改善, 能在相同墙钟预算下完成更多有效运算. 浮点算术量一样的两个实现也可能运行时间不同. 因此算法 FLOPs、设备峰值、实测吞吐与总时间是四个不同量.

为准确续写多种文本, 模型可能需要表示实体、关系、语法、代码状态与证明中的依赖. 这是 ``条件预测压力促进有用表示'' 的机制假说. 它不证明每一个低损失模型都形成正确因果世界模型: 有限资料可被记忆, 同一输出分布可以有不同内部实现, 数据中的错误规律也可能被准确拟合. 从生成成绩到一般智能或主观体验, 还隔着对象定义和证据要求, 不能由一个损失函数直接推出.

\originalsection{训练 FLOPs 的近似账本}\label{sec:progress-flops}
训练预算常用浮点运算次数计量, 简称 FLOPs. 本节采用一次标量乘法与一次加法各算一次的约定, 所以一个乘加大约算两次. 设矩阵 $W\in\R^{m\times d}$, 一批 $S$ 个 token 隐状态为 $X\in\R^{S\times d}$, 则前向 $XW^\top$ 约需 $2Smd$ 次运算. 反向计算输入梯度与权重梯度各有一个同级矩阵乘法, 合计前后向约为 $6Smd$.

如果网络中的多数参数都在这种稠密矩阵中, 每个有效 token 各使用一次, 把参数数 $N$ 与累计训练 token 数 $D$ 相乘, 得到
\begin{equation}\label{eq:compute-six}
C_{\rm train}\approx6ND.
\end{equation}
这里 $D$ 是处理次数, 不是互不重复的文本量. 一份长度 $T$ 的训练文档重复十轮就贡献约 $10T$ 个 token 运算, 但信息量没有因此自动乘十. Embedding 查表中的参数并不是每个 token 都作稠密乘法; 输出投影却会涉及整个词表. 因而 $N$ 的计数约定应写清, 不能把任何参数总数都无条件代入.

注意力矩阵另有依赖序列长度的成本. 一层宽度 $d$、长度 $T$ 的 $QK^\top$ 与 $AV$ 各需约 $2T^2d$ 次前向运算; 各自的反向再增加约两倍, 前后向共约 $12T^2d$. $L$ 层、累计 $D$ 个 token 若由等长文档构成, 这部分约为 $12LDTd$. 相比投影与前馈的 $O(LDd^2)$, 长上下文时它不能忽略. 稀疏路由模型还要区分总参数与每 token 激活参数; 重计算会增加 FLOPs; optimizer 更新、通信、数据读取与检验也有成本. 式 \eqref{eq:compute-six} 是稠密、矩阵运算占主导时的尺度近似.

\begin{example}
取 $N=10^9,D=2\times10^{10}$, 忽略额外注意力与系统成本. 若实测持续吞吐为 $10^{14}$ FLOPs/s, 估计核心训练时间. 若该吞吐实际只有设备峰值的 $1/4$, 能否仍用峰值作除数?
\end{example}
\begin{solution}
近似计算为 $6ND=1.2\times10^{20}$ FLOPs, 除以持续吞吐得 $1.2\times10^6$ 秒, 约 13.9 天. 应使用实际完成本任务运算的吞吐; 用四倍峰值作除数会低估约四倍. 检查点、评测与数据停顿还可能增加总时间.
\end{solution}
这种账本能解释为什么训练大模型需要同时优化算法与系统, 也说明论文中的同 FLOPs 比较和实际同时间比较并非相同实验. 本书的 NumPy 小模型只用于机制复算, 并未测量大型分布式训练的利用率.

\originalsection{规模损失与预算最优点}\label{sec:progress-scaling}
规模规律研究的是: 固定模型族、数据分布与训练协议, 最终损失怎样随模型大小和处理资料改变. Kaplan 等论文报告过近似幂律关系\cite{kaplan}. Hoffmann 等在 Chinchilla 研究中训练了 400 余个模型, 覆盖约 70M--16B 参数、5B--500B token, 并比较几种预算分配估计方法; 最后以 70B 参数、1.4T token 训练 Chinchilla, 与 280B 参数 Gopher 的训练预算比较\cite{chinchilla}. 这些是作者的实验范围与报告, 不是本项目复现的数据.

为了完整理解预算平衡, 采用一个待拟合模型
\begin{equation}\label{eq:scaling-model}
L(N,D)=L_\infty+aN^{-\alpha}+bD^{-\gamma},
\qquad a,b,\alpha,\gamma>0.
\end{equation}
$L_\infty$ 表示该模型形式假定的极限损失, 两项分别随参数与 token 增加而下降. 本节用 $\gamma$ 表示数据指数, 避免和后训练的 KL 系数 $\beta$ 混淆. 加法分解是经验模型假设; 它没有由概率链式法则自动推出. 对优化不足、重复数据或分布变化, 这个形式可能拟合不好.

给定数据点 $(N_i,D_i,L_i)$, 可通过最小化预测与观测损失的某种残差拟合五个常数. 因为 $L_\infty$ 与幂项之和需要先减再取对数, 不能简单把 $\log L$ 对 $\log N$ 回归就声称得到所有指数. 参数与 token 总是按一个固定比例同时变化时, 两类效应也难以分别识别. 较可信的设计要改变两个轴, 留出若干尺寸或预算检验预测, 对训练停止与学习率安排作控制, 再报告拟合不确定性与外推限制.

\begin{theorem}\label{thm:compute-optimal}
在式 \eqref{eq:scaling-model} 成立且约束为 $C=kND$、$k>0$ 的连续近似下, 唯一损失最小点为
\begin{align*}
N_*&=\left(\frac{\alpha a}{\gamma b}\right)^{1/(\alpha+\gamma)}
       \left(\frac Ck\right)^{\gamma/(\alpha+\gamma)},\\
D_*&=\left(\frac{\gamma b}{\alpha a}\right)^{1/(\alpha+\gamma)}
       \left(\frac Ck\right)^{\alpha/(\alpha+\gamma)}.
\end{align*}
在该点, $\alpha aN_*^{-\alpha}=\gamma bD_*^{-\gamma}$.
\end{theorem}
\begin{proof}
约束使 $D=C/(kN)$, 因而可变部分为
$f(N)=aN^{-\alpha}+b(kN/C)^\gamma$, $N>0$. 求导并乘以正数 $N$, 得
\[
Nf'(N)=-\alpha aN^{-\alpha}+\gamma b(kN/C)^\gamma.
\]
第一项绝对值随 $N$ 严格下降, 第二项严格上升. 当 $N$ 从零增到无穷时, 导数符号只由负变正一次. 解平衡式,
\[
N^{\alpha+\gamma}=\frac{\alpha a}{\gamma b}(C/k)^\gamma,
\]
得 $N_*$, 再用约束得 $D_*$. $f$ 在此点之前下降、之后上升, 所以这是唯一全局最小点. 最后的等式就是驻点条件.
\end{proof}
边际平衡并不是让两项误差总相等. 若指数不同, 平衡的是指数乘误差. 用对数预算理解更直接: 增加 $\log N$ 会以速度 $\alpha aN^{-\alpha}$ 降低第一项, 同时固定计算迫使 $\log D$ 等量下降, 使第二项以速度 $\gamma bD^{-\gamma}$ 上升. 最优点使两者相等.

若 $\alpha=\gamma$, 参数与数据都随 $C^{1/2}$ 增长, 但前面的常数仍依赖 $a/b$. 若 $\alpha=1/2,\gamma=1$, 参数随 $C^{2/3}$、数据随 $C^{1/3}$ 增长. 计算增加八倍时, 两者分别增加四倍与两倍. 最优超额损失的两项都随 $C^{-\alpha\gamma/(\alpha+\gamma)}$ 缩减, 可以代入上面的解直接验证. 这一推演是本书对明确假设的数学分析, 不代表任何现实模型自动满足同样指数.

Chinchilla 第 3 节中不同估计路线给出的最优指数并非逐位相同, 其中固定模型训练曲线方法报告接近 0.50/0.50, 等 FLOPs 方法报告约 0.49/0.51\cite{chinchilla}. 因而应把 ``参数与数据协同增长'' 当作该实验范围内的经验发现, 而非把某一个 token/参数比奉为所有训练的固定律. 本项目没有多尺寸多预算语言训练网格, 不能用一条小模型训练曲线检验这些幂律.

\originalsection{部署成本改变预算选择}\label{sec:progress-deploy}
若训练后要服务大量请求, 目标不能只含一次训练成本. 对稠密网络, 每生成一个 token 的矩阵前向主项近似为 $2N$ FLOPs, 还需计入注意力、输出、缓存读写与检索工具. 设生命周期生成 token 数为 $M$, 用简化总账本
\[
C_{\rm total}=kND+qNM,
\]
其中训练常取 $k\approx6$, 前向矩阵主项常取 $q\approx2$. $M$ 是部署需求, 不是训练资料. 将固定预算 $C_0$ 解为 $D=C_0/(kN)-qM/k$, 需要 $0<N<C_0/(qM)$, 才有正训练 token 预算.

令 $h=qM/k$. 代入损失模型并对 $N$ 求导, 注意 $D'(N)=-C_0/(kN^2)$, 得
\[
f'(N)=-\alpha aN^{-\alpha-1}
+\gamma bD^{-\gamma-1}\frac{C_0}{kN^2}.
\]
因此驻点满足
\[
\alpha aN^{-\alpha}=\gamma bD^{-\gamma}(1+h/D).
\]
与仅训练的条件相比, 右边多出 $1+h/D$: 模型变大既挤占训练 token, 又增加每次部署成本. 这个模型中, 大量服务需求倾向选择更小模型、训练更多 token. 例如在训练最优点之后略增大 $N$, 仅训练目标已上升; 总成本约束下 $D$ 比仅训练情形还小, 正的上升项更大, 因而部署最优点不会由这个效应移到更大的 $N$ 一侧. 现实成本还受延迟、批大小、量化与设备内存约束, 不宜仅用 FLOPs 决定部署.

缓存把重复计算换成存储, 所以单请求更快也可能减少设备能并发容纳的请求数. 输出很长时, $M$ 增长与缓存长度增长同时发生. 更强的小模型、蒸馏、量化或更准确的早停都有可能降低成本, 但它们也改变质量条件. 比较两个产品应明确质量、延迟与总吞吐的约束, 不能只报告 ``参数更少''.

\originalsection{资料质量改变同一 token 预算}\label{sec:progress-data}
设训练混合分布为 $P=\sum_jw_jP_j$, $w_j\ge0,\sum_jw_j=1$. 对固定模型 $q$, 期望交叉熵为 $\sum_jw_jH(P_j,q)$. 改变权重即改变学习目标. 如果评价分布为 $Q$, 降低 $P$ 上损失不一定降低 $Q$ 上损失. 例如一个模型容量只允许输出固定字母分布, 训练将 90\% 权重放在 ``a'' 上, 而测试把 90\% 放在 ``b'' 上, 训练频率拟合得更准也不能解决任务分布错位.

重复资料可以帮助优化收敛, 也会使某些文本被过度加权. 若一份文档复制十遍, 最大似然目标中它的权重就增大十倍; 算力增加并没有创造十份独立证据. 近重复的模板、网页导航与乱码也可能占去大量 token. 质量过滤尝试把预算用于更有用的模式, 但过滤过强可能删除小众知识、方言或困难样本, 将偏好误当正确性.

FineWeb 论文的 v2 版本通过小模型训练来比较抽取、过滤与去重的选项, 其第 3 节含去重范围的对照, 第 4 节含教育质量过滤与下游评价\cite{fineweb}. 它的价值不只是提供资料名称, 而是把 ``清理更好'' 变成固定训练设置下可观察的消融. 其中第 3.4 节的具体反例是: 按时间跨全部抓取全局去重, 对旧抓取可删去约 90\% 资料, 留下的资料反而在其训练评价中较差; 逐抓取独立去重得到更好的平均成绩. 这说明更激进地删重复不必然更好, 去重还改变了保留内容的分布. 因而既要报告过滤规则也要报告可用 token 总数. 这些作者结果受模型尺寸、语言、数据混合与基准集合限制, 不证明同一个过滤器在全部模型与语言上最佳.

评价污染尤其容易混淆. 若训练中包含测试题及答案, 模型成绩可能反映重现, 而非迁移. 去除字符串完全相同的题面不足以处理改写、答案解析或由教师生成的近似题. 因而资料质量研究同时需要检查测试重合、来源划分与新任务保留点. 合成资料也需要这个审查: 教师能生成流畅文本, 但错误、模式重复和任务窄化可能被学生一起吸收.

\originalsection{阈值、涌现与行为边界}\label{sec:progress-emergence}
评价常把连续概率压成离散得分. 精确匹配只记录完整答案对错, 部分得分则保留一些结构. 若一个任务要求 $m$ 个步骤全正确, 简化地假设各步独立且正确率相同为 $p$, 则总成功率为 $p^m$. 当 $m=20$, $p=0.7$ 时约为 0.000798, $p=0.9$ 时约为 0.121577. 每步只是连续改善, 完整成功却可能从测不到变成明显可见.

这个变化还是连续函数, 没有数学断点. 若每个规模只测几十题, 小概率成功可能暂未出现; 若绘图仅标 ``是否超过某阈值'', 连续改善就成为阶跃标签. Schaeffer 等的涌现指标论文讨论了指标选择和统计分辨率能怎样影响这类图像\cite{emergence}. 这提醒我们同时观察概率分数、部分得分、精确正确率和置信区间, 不意味着所有结构变化都由指标造成.

步骤独立是上述算例的条件. 若错误共享同一误解, $p^m$ 无效; 若不同题有不同 $p$, 先平均 $p$ 再幂运算也无效. 例如一半题可完美解决、另一半完全不能解决, 平均单次正确率为 $1/2$. 在同一题反复独立采样也不会让后半题变好, 整套覆盖率始终是 $1/2$, 不能套 $1-(1-1/2)^N$ 宣称趋于 1. 这类异质性会在下一章的预算曲线中再次出现.

``能力出现'' 至少可以表示三种不同主张: 在既定任务上超过可观测阈值; 系统能完成以前未完成的任务; 网络内部形成新的计算机制. 前两者可以由行为证据支持, 第三者需要机制研究. 即使内部机制改变, 是否与人类智能同构、是否具有意识, 仍不是由一个 benchmark 分数决定的事实. 本书采用能检验的条件预测、迁移与任务成绩, 把更广的哲学判断留在其证据范围之外.

\originalsection{练习与解答}\label{sec:progress-practice}
\begin{exercise}
在模型 $L=L_\infty+aN^{-1/2}+bD^{-1}$ 中, 某一点的两项超额损失恰好相等. 在固定 $ND$ 下它是否最优? 应怎样调整?
\end{exercise}
\begin{solution}
最优要求 $(1/2)aN^{-1/2}=bD^{-1}$, 不是两项相等. 若两项均为 $e>0$, 增加 $\log N$ 的导数为 $-e/2+e=e/2>0$. 应减小 $N$、增加 $D$, 直到两类边际收益平衡. 只比较误差数值而忽略指数会选错方向.
\end{solution}
\begin{exercise}
某团队把同一训练语料重复四轮, 声称 ``数据量增加四倍, 所以验证了数据规模规律''. 哪个量确实增加? 需要怎样设计才能检验新资料的效应?
\end{exercise}
\begin{solution}
处理 token 数与核心训练计算增加四倍, 不同文本数量没有增加. 可设重复旧资料与加入同分布新资料两个组, 保持总处理 token、模型和优化基本一致, 在独立保留资料上比较. 还要检查两组是否已收敛, 不能把额外训练步数和新信息混为一谈.
\end{solution}
\begin{exercise}
模型甲单次正确率较高, 模型乙允许多生成 100 次并使用精确 oracle 选择后成绩更高. 这支持哪一项结论? 怎样比较实际部署系统?
\end{exercise}
\begin{solution}
支持乙在该生成与 oracle 协议下有较好的系统结果, 不能单独归因为权重更强. 应报告单次成绩、候选覆盖、真实选择器成绩、总 token 与时间, 并比较相同成本的预算曲线. 如果部署不能获得 oracle, 还要测替代验证器的误拒与错误接受.
\end{solution}
\label{ch:end:09}
